\documentclass[12pt, a4paper]{article}

\usepackage{amsmath}
\usepackage{esint}
\usepackage{derivative}

\newcommand{\manualindent}[1]{
    \section{#1}
    \noindent\hspace*{1.5em} % Adjust this space as needed
}

% enable and disable indentation by inserting the 
\makeatletter
\NewCommandCopy{\indentation}{\@afterindentfalse} %create a deep copy of \@afterindentfalse command as \indentation

\newcommand{\indenton}{\let\@afterindentfalse\@afterindenttrue} % set indent to on
\newcommand{\indentoff}{\let\@afterindentfalse\indentation} % set indent to off
\makeatother
% \makeatletter
% \let\@afterindentfalse\indentate
% \makeatother

\bibliographystyle{plain}




% \setlrmarginsandblock{3cm}{3cm}{*}
% \setulmarginsandblock{3cm}{3cm}{*}

\begin{document}

% prints hello world
Hello, World!
% 
% \ShowCommand{\@afterindenttrue}

\vspace{10cm} %create 10cm of vertical space


\section{Most common \LaTeX~commands}


\begin{verbatim}
    \def

    Defines a new macro.
    Syntax:
    \def\<macro>{definition}
    Example:
    \def\square#1{#1^2} % Defines a macro square which squares its input.
    When compared with \newcommand, note that \def is more primitive and doesn’t check if a command already exists—overwriting it blindly.
    \newcommand
    
    A safer alternative to defining new macros than using plain TeX’s \def.
    Syntax:
    \newcommand{\<macro>}[num_args]{definition}
    Example:
    \newcommand{\myname}{John Doe}
    \renewcommand
    
    Used if you need to redefine an existing macro that was previously defined by someone else.
    Syntax:
    \renewcommand{\<macro>}[num_args]{definition}
    \expandafter
    
    Used for advanced TeX programming where controlling expansion timing matters.
    Example use case: Changing expansion order of macros before execution.
    \csname ... \endcsname
    
    Allows dynamic generation of control sequences.
    Useful when you want programmatic access to macros whose names are not known beforehand but built from other variables.
    `\if ... \fi Constructs These conditionals allow basic logic within LaTeX documents.
    
    **\relax A no-op (does nothing). Often used as a placeholder or delimiter marker in complex macros.
    
    Environment Handling (begin{} & end{}) Essential for creating structured blocks like equations (equation*) or itemized lists (itemize).
\end{verbatim}
    

% math mode
\hspace{-3cm} 1. This will start math mode and line break examples:

Math mode: \hfill \break
\hspace{5cm} The equation is \(E = mc^2\) % inline math, don't know why \hspace does not work.
\[
    E=mc^2 % display math
\]

Math mode: \\
\hspace{2cm}The equation is $E = mc^2\$$ % inline math not recommended with $
\[
    E=mc^2 % display math
\]

Math mode: \newline
The equation\\* is \(E = mc^2\) % inline math
\[E=mc^2\]

% labels and references
We know from equation~\ref{eq:emc} that $E=mc^2$.\\ % with tilde means a space is added, but if there's a line break, the equation + number will be on the next line
We know from equation \ref{eq:emc} \dotfill that $E=mc^2$. % without tilde means no space is added

\begin{equation}
    E=mc^2
    \label{eq:emc}
\end{equation}

We know from Equation~\ref{eqn:emc}:
\begin{equation}
    E = mc^2
    \label{eqn:emc}
\end{equation}

\section{This will start a new section with a number and title}

This is the \textbf{bold} text.

\section[short title]{This will start a new section with a number and title, but the title in the table of contents will be `short title'} % the short title is an optional argument

This is the \textit{italic} text.

\section*{This will start a new section without a number and title}

This is the \underline{underlined} text.

\section[matrices and integrals]{Matrices and integrals expressions}

First, import \verb|amsmath| \verb1package1. % the \verb| | can use different delimiters, such as 1 or <>, it doesn't matter as long as they match and doesn't appear through the content. It does not use curly braces, because curly braces are interpreted are reserved for latex grouping


\subsection{Matrices 1} % depends on amsmath package 
Here is a matrix:
\[
    A= \begin{pmatrix}
        a & b \\
        c & d \\
        1 & 2 \\
    \end{pmatrix}
\]

\subsection{Integral 1}

\[
    \int_0^\infty x^2 \,dx > \int_{a}^{b} x^2*y \,dx\,dy
\]

\subsection{Multiple Integrals}

\indent This is a multiple integral that should be indented % this line is not being indented even after the \indent command

\[
    \iiiint_V f(x,y,z,w) \,dx\,dy\,dz\,dw
\]

This is a cyclic multiple integral (does not exist and does not make sense)

\[
    \oint_C f(x,y,z) \,dx\,dy\,dz
\]\linebreak

This is another multiple integral with fraction:

\[
\iiiint_V \frac{f(x,y,z,w)}{2}x^2y^2+z^2+w^2 \,dx\,dy\,dz\,dw
\]

\indenton

\section{Derivatives}

\subsection{Partial Derivatives}
to write partial derivatives:
% Partial derivative symbol in LaTeX
\[
   \frac{\partial f}{\partial x_i}
\]

\subsection{using derivatives package}
\subsubsection{By using package derivatives}
\[
\pdv{f}{x_2, y^3, z} = \pdv{e^{z}\sin(x)}{x} + \pdv{\frac{e^{y}}{\cos(y)}}{y} + \pdv{e^{z}\cot(z)}{z}
\]

\subsubsection{Evaluate at a specific point}

Some example equations
\begin{equation}
    \pdv{U}{x}_p = \pdv{U}{x}\bigg|_p = \pdv{U}{x}\big|_{p, Y}\quad = \pdv{U}{x}\Bigg|_{(p, Y)}\quad
    \label{eq:partial_derivative}
\end{equation}


\section{section Without Indentation}

text2 should be indented

\subsection{This is an unindented subsection}

tex should be indented

\indentoff
\subsection{Subsection With Indentation}

This should be unindented

\phantom{hidhsi}

sdasdsda

\section{Citation}

This is a citation \cite{tcc_vitorhp}.

\bibliography{refs}
% Next in line: multiple types of document classes, packages, and environments


\end{document}
